\documentclass[11pt,a4paper]{article}

\usepackage[T1]{fontenc}
\usepackage{lmodern}
\usepackage[french]{babel}
\usepackage[margin=2.3cm]{geometry}
\usepackage{array,booktabs,tabularx}
\usepackage[table]{xcolor}
\usepackage{fancyvrb}
\usepackage[hidelinks,linktoc=all]{hyperref}

\definecolor{encre}{HTML}{173451}
\definecolor{psbar}{HTML}{242933}
\definecolor{psbody}{HTML}{012456}
\definecolor{pstext}{HTML}{F1F5FB}

\setlength{\parindent}{0pt}
\setlength{\parskip}{5pt}
\renewcommand{\arraystretch}{1.2}
\setcounter{tocdepth}{3}
\setcounter{secnumdepth}{3}

\newcommand{\etudiant}{Kamel Mahdi Gherbi}
\newcommand{\code}[1]{\texttt{#1}}
\newcommand{\flux}{\ensuremath{\rightarrow}}
\newcommand{\preuve}[2]{%
  \href{run:../preuves/#1}{\textcolor{encre}{\underline{#2}}}}
\newcommand{\manifeste}[2]{%
  \href{run:../data/manifestes/#1}{\textcolor{encre}{\underline{#2}}}}

\newsavebox{\psoutput}
\newcommand{\afficherpowershell}{%
  \par\smallskip
  \begingroup
  \setlength{\fboxsep}{10pt}%
  \noindent\colorbox{psbar}{%
    \makebox[\dimexpr\linewidth-20pt\relax][l]{%
      \color{pstext}\small\sffamily\bfseries PowerShell
      \hfill\mdseries $-$\quad$\times$}}%
  \par\nointerlineskip
  \noindent\colorbox{psbody}{\usebox{\psoutput}}\par
  \endgroup
}

\begin{document}

\begin{titlepage}
  {\small\bfseries UNIVERSITÉ PARIS-SACLAY \quad M1 MIAGE \quad UE DSID\par}
  \vspace{2.2cm}
  {\Huge\bfseries\color{encre} Projet Big Data\par}
  \medskip
  {\Large Chaînes d'ingestion et contrôles Bronze--Silver\par}
  \bigskip
  \etudiant \hfill Année universitaire 2026--2027
  \vfill

  \textbf{Avancement du projet}

  Le flux~1, consacré à l'API des prix des carburants, a produit un lot
  brut vérifié dans MinIO. Le chargement PostgreSQL Bronze et les
  contrôles de la vue Silver ont été validés sur les 9\,834 stations
  de ce lot. Les cinq connexions des trois flux sont actives et la
  vue d'ensemble Airbyte indique un état sain. Le contrôle Bronze
  du flux CSV a confirmé 101 codes département distincts.
  Le contrôle Silver du flux CSV a confirmé le même volume et la
  présence des champs de traçabilité sur les 101 lignes.
  Le contrôle Bronze confirme l'arrivée des quatre nouvelles fiches CRM.
  Le rattachement de l'opportunité au compte est vérifié dans Bronze.
  Les contrôles de volume, de clés et de traçabilité des trois vues Silver du CRM sont validés.
\end{titlepage}

\tableofcontents
\clearpage

\section{Architecture}
\label{sec:architecture}

\subsection{Périmètre et environnement local}

Le projet a été exécuté sur un poste Windows avec Docker
Desktop~4.92.0, Docker Engine~29.8.0 et Docker Compose~v5.5.1.
Docker Desktop exécute les conteneurs en \code{linux/amd64}.
Docker Compose démarre les services du kit sur le réseau
\path{dsid-bigdata-network}. Airbyte est déployé séparément dans un
cluster Kubernetes local \code{kind}, créé avec \code{abctl v0.30.4}.
Le chart Helm et l'application Airbyte sont tous deux en version~2.5.0.
Les sorties d'inventaire figurent à la section~\ref{sec:docker-airbyte}.

{\small
\begin{tabularx}{\linewidth}{@{}>{\bfseries}p{2.9cm}p{3.3cm}X@{}}
\toprule
Composant & Version ou référence & Rôle et accès local \\
\midrule
Airbyte & 2.5.0 &
Exécute les synchronisations. Interface Web sur le port~8000. \\
MinIO & Image construite par le kit &
Stockage objet compatible S3. API sur~9000, console sur~9001.
La référence de construction et la version du binaire actif
figurent à la section~\ref{sec:docker-airbyte}. \\
PostgreSQL & 18.6 &
Base \code{warehouse} sur~5432. \code{bronze} reçoit les données
chargées et \code{silver} expose les vues SQL standardisées. \\
MariaDB & 11.8.9 &
Base opérationnelle \code{espocrm} sur~3306. \\
EspoCRM et daemon & 10.0.8 &
Saisie des fiches et exécution des tâches du CRM. Interface Web sur~8081. \\
Adminer & 6.1.0 &
Interface de consultation des bases et d'exécution des contrôles SQL
sur~8082. \\
\bottomrule
\end{tabularx}
}

\subsection{Réseau et types de connexion}

Le navigateur du poste accède aux interfaces Web via \code{localhost}.
Les connecteurs Airbyte s'exécutent dans le cluster \code{kind}, sur un
réseau distinct de Compose. Pour joindre les ports publiés sur la machine
hôte, ils utilisent l'adresse \path{host.docker.internal}.

{\small
\begin{tabularx}{\linewidth}{@{}>{\bfseries}p{3.2cm}X@{}}
\toprule
Liaison & Protocole et configuration utilisés \\
\midrule
Sources publiques &
Requête HTTPS GET sans clé pour l'API carburants.
Lecture HTTPS du CSV Insee avec le connecteur File. \\
MinIO &
Connexion S3 compatible via \url{http://host.docker.internal:9000}.
Le bucket \code{dsid-raw} est privé. Les clés d'accès sont configurées
localement dans Airbyte. La lecture a été validée avec Source S3~4.16.4. \\
PostgreSQL &
Connecteur Postgres~3.0.24, hôte \code{host.docker.internal}, port~5432,
base \code{warehouse}, schéma cible \code{bronze}. \\
MariaDB &
Connecteur MySQL, hôte \code{host.docker.internal}, port~3306,
base \code{espocrm}. Pour la reprise, le kit crée le compte dédié
\code{airbyte} avec un droit de lecture \code{SELECT}. \\
\bottomrule
\end{tabularx}
}

\subsection{Parcours des données et synchronisations}

{\small
\begin{tabularx}{\linewidth}{@{}>{\bfseries}p{3cm}X@{}}
\toprule
Flux & Parcours des données \\
\midrule
1. API carburants &
API publique \flux\ Airbyte (connecteur Builder) \flux\ MinIO RAW
\flux\ Airbyte Source S3 \flux\ PostgreSQL Bronze \\
2. Fichier CSV &
CSV Insee \flux\ Airbyte File Source \flux\ MinIO RAW
\flux\ Airbyte Source S3 \flux\ PostgreSQL Bronze \\
3. EspoCRM &
EspoCRM \flux\ MariaDB \flux\ Airbyte MySQL Source
\flux\ PostgreSQL Bronze \\
\midrule
Couche Silver &
PostgreSQL Bronze \flux\ scripts SQL par source
\flux\ PostgreSQL Silver \\
\bottomrule
\end{tabularx}
}

Les cinq connexions Airbyte utilisent un déclenchement manuel. L'API
et le CSV ont chacun deux connexions, l'une pour déposer les données
dans MinIO et l'autre pour les charger en Bronze. Le CRM utilise une
connexion directe de MariaDB à PostgreSQL. Pour l'API, le mode
\emph{Full Refresh -- Overwrite} relit le périmètre sélectionné
et maintient un état courant. Le manifeste identifie le lot d'origine
par sa clé objet et son empreinte SHA-256.

\preuve{airbyte/connexions_vue_ensemble.png}{Vue des cinq connexions Airbyte}.
Elle montre les sources, les destinations et le déclenchement manuel.

\section{Sources}

\subsection{Flux 1 --- API publique des prix des carburants}

\subsubsection{Choix de la source et contrat de données}

Le jeu \emph{Prix des carburants en France, flux instantané v2} est
publié par la DGCCRF. L'API Explore v2.1 expose les données en JSON
par requête HTTPS GET, sans clé. Le jeu est publié sous Licence
Ouverte 2.0 avec attribution de la source.

Chaque objet du tableau \code{results} représente une station dans
l'état observé pendant la collecte. Le champ \code{id} identifie
la station. Le lot identifie une collecte et permet de distinguer
les observations successives d'une même station.

{\small
\begin{tabularx}{\linewidth}{@{}>{\bfseries}p{2.9cm}X@{}}
\toprule
Groupe & Exemples de champs et de types observés \\
\midrule
Identification &
\code{id} numérique dans le JSON, code postal \code{cp}, ville et
adresse en texte. \\
Prix &
\code{gazole\_prix} et prix des autres carburants numériques.
Les champs \code{*\_maj} sont des chaînes horodatées. \\
Disponibilité &
Tableaux de carburants disponibles ou indisponibles et champs
décrivant les ruptures publiées. \\
\bottomrule
\end{tabularx}
}

\href{https://data.economie.gouv.fr/explore/dataset/prix-des-carburants-en-france-flux-instantane-v2/api/}{Console API du ministère}
\quad
\href{https://www.data.gouv.fr/datasets/prix-des-carburants-en-france-flux-instantane-v2-amelioree}{Jeu de données et licence}

\subsubsection{Test du GET dans PowerShell}

Le filtre \verb|where=startswith(cp, '75')| a limité le test aux
stations dont le code postal commence par 75. Le tri explicite
\verb|order_by=id| facilitait la comparaison des pages. La taille
de~20 a permis de contrôler deux changements de page et l'arrêt sur
une dernière page partielle.

\begin{lrbox}{\psoutput}
\begin{minipage}{\dimexpr\linewidth-20pt\relax}
\color{pstext}\small
\begin{Verbatim}
HTTP 200 | offset=0  | count=20 | total_count=47
HTTP 200 | offset=20 | count=20 | total_count=47
HTTP 200 | offset=40 | count=7  | total_count=47

Stations recues         : 47 / 47
Identifiants distincts  : 47
Identifiants absents    : 0
Doublons                : 0
Stations hors perimetre : 0
Validation              : True
\end{Verbatim}
\end{minipage}
\end{lrbox}
\afficherpowershell

\preuve{api/api_01_carburants_test_terminal.png}{Capture du test PowerShell}.
Elle montre les réponses HTTP, les volumes par page et les contrôles
des identifiants et du périmètre.

\subsubsection{Configuration et validation dans Airbyte Builder}

Le connecteur utilise la
\href{https://data.economie.gouv.fr/api/explore/v2.1/catalog/datasets/prix-des-carburants-en-france-flux-instantane-v2/records}{route nationale \code{/records}}.

{\small
\begin{tabularx}{\linewidth}{@{}>{\bfseries}p{3.1cm}X@{}}
\toprule
Réglage & Choix et justification \\
\midrule
Méthode et accès &
GET sans authentification pour interroger la route publique. \\
Flux &
\code{prix\_carburants\_france} dans le YAML Builder versionné.
Le flux S3 qui relit le lot s'appelle \code{prix\_carburants}. \\
Sélecteur &
\code{results}, tableau qui contient les stations. \\
Clé &
\code{id}, identifiant métier de la station. \\
Pagination &
\emph{Offset Increment}. \code{limit} indique la taille de page
et \code{offset} indique le nombre d'enregistrements à sauter. \\
\bottomrule
\end{tabularx}
}

Le test Builder sur Paris a vérifié l'extraction de \code{results}
et le passage entre les trois pages. Pour la collecte nationale,
le filtre et le tri explicite du test ont été retirés. L'absence de
\code{select} et de transformation Builder conserve les champs et
les valeurs transmis par la source.

La taille de page est passée de~20 à~100 pour réduire le nombre
de requêtes. L'aperçu national a confirmé les décalages 100 et 200.
La documentation de la route fixe un plafond de 10\,000 pour
\code{offset + limit}. La complétude se vérifie en comparant le
nombre collecté au \code{total\_count} annoncé pendant le même job.
Les relevés du projet ne permettent pas d'établir le quota de débit
applicable à cette route.

\preuve{api/api_03_carburants_airbyte_builder.png}{Capture du test Builder}.
Elle montre la configuration et l'aperçu du flux de test.
\preuve{api/api_04_carburants_pagination_nationale.png}{Pagination du flux national}.
Cette capture montre le réglage à 100 enregistrements par page et
l'aperçu Builder.
\quad
\href{https://help.opendatasoft.com/apis/ods-explore-v2/}{Documentation de la pagination}

\subsection{Flux 2 --- Fichier CSV}

\subsubsection{Choix du fichier et contrat de données}

Le jeu retenu est le fichier des départements du Code officiel
géographique 2026 publié par l'Insee. Une ligne décrit un département,
identifié par le code \code{DEP}. Ce code est traité comme du texte
afin de préserver les zéros initiaux et les codes alphanumériques.
Le CSV comporte un en-tête, un séparateur virgule et un encodage UTF-8.
Le fichier original contient 101 enregistrements et 101 codes
distincts. L'Insee le diffuse sous Licence Ouverte.
Une copie du fichier est fournie sous \path{data/v_departement_2026.csv}.

\href{https://www.data.gouv.fr/datasets/code-officiel-geographique-1}{Fiche du jeu de données et licence}

\href{https://www.insee.fr/fr/statistiques/fichier/8740222/v_departement_2026.csv}{Fichier CSV source de l'Insee}

\subsubsection{Contrôle du fichier brut}

Le connecteur File lit le CSV via HTTPS. Airbyte écrit la sortie datée
dans le bucket \code{dsid-raw}, sous \path{raw/csv/departements}.
La synchronisation du 13 novembre 2026 a chargé 101 enregistrements
dans MinIO. Après téléchargement de la copie archivée de l'original,
le contrôle PowerShell a confirmé que son SHA-256 correspond à celui
du fichier officiel. Le CSV généré par Airbyte a une empreinte
différente. Le manifeste consigne les empreintes, la clé de l'objet
généré et le nombre d'enregistrements chargés.

\begin{itemize}
  \item \preuve{csv/csv_01_departements_sync_minio.png}{Synchronisation CSV vers MinIO}
        \enspace Airbyte indique 101 enregistrements chargés.
  \item \preuve{csv/csv_02_departements_objet_minio.png}{Objet CSV dans MinIO}
        \enspace La console affiche l'objet daté produit par Airbyte.
  \item \preuve{csv/csv_03_departements_sha256.txt}{Contrôle SHA-256 de l'original archivé}
        \enspace Le résultat PowerShell \code{True} confirme l'égalité des empreintes.
  \item \manifeste{csv_departements_2026.json}{Manifeste du lot CSV}.
\end{itemize}

\subsection{Flux 3 --- EspoCRM et MariaDB}

\subsubsection{Périmètre de la source}

La source MySQL \code{mariadb\_espocrm} lit la base MariaDB
\code{espocrm} et réplique les flux \code{account}, \code{contact} et
\code{opportunity} vers PostgreSQL Bronze. Chaque ligne représente une
fiche de son entité et la clé métier est le champ \code{id}.
Les colonnes \path{contact.account_id} et \path{opportunity.account_id}
portent les rattachements au compte vérifiés dans ce scénario.

Une capture de configuration affiche l'utilisateur \code{espocrm}.
Le compte utilisé par le job réussi n'est pas établi par les preuves
du dépôt. La procédure de reprise utilise le compte MariaDB dédié
\code{airbyte} du kit et vérifie son droit \code{SELECT} avant le test
de connexion.

\section{Bronze}

\subsection{Flux 1 --- De MinIO à PostgreSQL}

\subsubsection{Lot brut contrôlé dans MinIO}

Le job Airbyte~\code{5} a chargé 9\,834 enregistrements. MinIO affiche
un fichier JSONL de 25,4~MiB dans le bucket privé \code{dsid-raw}.

\begin{quote}
\small
\path{raw/api/prix_carburants/extraction_date=2026-11-13/1794580622068_0.jsonl}
\end{quote}

Le fichier téléchargé contient 9\,834 lignes, soit le volume chargé
affiché dans Airbyte. Le \code{total\_count} national du job n'est pas
consigné. La concordance des volumes établit la cohérence du lot
collecté et chargé.

\begin{lrbox}{\psoutput}
\begin{minipage}{\dimexpr\linewidth-20pt\relax}
\color{pstext}\small
\begin{Verbatim}[fontsize=\footnotesize]
Lignes JSONL : 9834
Job Airbyte lu dans _airbyte_meta.sync_id : 5
SHA-256 :
AD91F416382A02FCEA107849390F3AF73449B5BF95AB29027CB9D9870214C80C
\end{Verbatim}
\end{minipage}
\end{lrbox}
\afficherpowershell

\preuve{api/api_05_carburants_sync_minio.png}{Synchronisation API vers MinIO}. La chronologie Airbyte montre
une exécution réussie avec 9\,834 enregistrements chargés.

\preuve{api/api_06_carburants_objet_minio.png}{Objet JSONL dans MinIO}. La console affiche la clé datée
et la taille de 25,4~MiB dans le bucket \code{dsid-raw}.

\manifeste{api_prix_carburants_1794580622068_0.json}{Manifeste du lot}.
Il associe le job, la clé MinIO, le volume et l'empreinte du fichier.

\subsubsection{Lecture S3 et configuration PostgreSQL}

La source \code{minio\_raw\_prix\_carburants} lit le bucket
\code{dsid-raw} en JSONL avec le motif suivant.

\begin{quote}
\small
\path{raw/api/prix_carburants/extraction_date=*/*.jsonl}
\end{quote}

Les caractères génériques acceptent les dates et les noms de fichiers
des prochains lots. Le manifeste conserve la clé exacte pour rejouer
une collecte précise.

La connexion \code{minio\_raw} \flux\ \code{bronze} utilise le schéma
de destination \code{bronze} et le préfixe de flux \code{api\_}.
Le mode \emph{Full Refresh -- Overwrite} recharge la table à partir
des objets sélectionnés par le motif S3. Le déclenchement manuel a
permis de vérifier le premier chargement.

\subsubsection{Incident de chargement et correction}

Les contrôles de connexion MinIO et PostgreSQL ont permis de retrouver
l'objet et d'atteindre le schéma cible. La réplication échouait ensuite
à la création de la table. Le journal de la destination indiquait
\verb|column "_airbyte_raw_id" specified more than once|.

Le JSONL contient les métadonnées du job API vers MinIO. Le connecteur
PostgreSQL génère ses propres colonnes techniques au chargement.
Pour éviter les collisions, quatre champs entrants ont été décochés
dans le schéma du flux \code{prix\_carburants}.

\begin{itemize}
  \item Champs décochés : \path{_airbyte_raw_id},
        \path{_airbyte_extracted_at}, \path{_airbyte_meta} et
        \path{_airbyte_generation_id}.
  \item Champs conservés : \path{_airbyte_data},
        \path{_ab_source_file_url} et
        \path{_ab_source_file_last_modified}.
\end{itemize}

La synchronisation a réussi après cette correction. Le JSONL dans
MinIO conserve les métadonnées du job~5. PostgreSQL contient les
métadonnées générées lors de son propre chargement et les deux champs
de provenance du fichier.

Les contrôles PostgreSQL établissent le résultat du transfert
MinIO vers Bronze. Le numéro~5 identifie le job API vers MinIO.
L'identifiant du job MinIO vers PostgreSQL ne figure pas dans les
preuves disponibles.

\subsubsection{Structure observée de la table Bronze}

La table \code{bronze.api\_prix\_carburants} contient sept colonnes.

{\small
\begin{tabularx}{\linewidth}{@{}l X@{}}
\toprule
Colonne & Type PostgreSQL \\
\midrule
\code{\_airbyte\_raw\_id} & \code{character varying} \\
\code{\_airbyte\_extracted\_at} & \code{timestamp with time zone} \\
\code{\_airbyte\_meta} & \code{jsonb} \\
\code{\_airbyte\_generation\_id} & \code{bigint} \\
\code{\_airbyte\_data} & \code{jsonb} \\
\code{\_ab\_source\_file\_url} & \code{character varying} \\
\code{\_ab\_source\_file\_last\_modified} & \code{character varying} \\
\bottomrule
\end{tabularx}
}

L'inventaire de \verb|_airbyte_data| a identifié 47 champs métier.
Leurs présences varient selon les stations. L'identifiant \code{id}
est accessible dans cet objet JSON.

\preuve{api/api_07_carburants_structure_bronze.png}{Structure de la table Bronze}. La requête affiche les sept
colonnes, les types PostgreSQL et les champs de provenance.

\subsection{Flux 2 --- Chargement du CSV}

La connexion \emph{minio\_raw vers bronze}, identifiée par la source
\code{minio\_raw\_departements}, figure parmi les cinq connexions
saines. La configuration documentée cible \code{bronze.csv\_departements}.
Le motif S3 retient les sorties CSV datées et exclut le dossier
\code{original/}. Le contrôle SQL dans la base \code{warehouse}
a confirmé 101 lignes dans \path{bronze.csv_departements}. Ce volume
correspond au nombre d'enregistrements du fichier officiel et au
nombre chargé par Airbyte vers MinIO. Le test compte 101 codes
département distincts, sans valeur absente ou vide et sans doublon.
Les codes \code{01}, \code{2A} et \code{2B} apparaissent chacun une fois,
ce qui confirme la conservation du zéro initial et des codes
alphanumériques vérifiés.

\preuve{csv/csv_04_departements_controle_bronze.png}{Requête et résultat du contrôle Bronze du CSV}.

\subsection{Flux 3 --- Chargement du CRM}

La connexion \emph{espocrm vers bronze} apparaît saine dans Airbyte.
L'inventaire SQL de la base \code{warehouse} confirme trois tables dans
le schéma \code{bronze}. Les volumes observés correspondent au relevé
initial de la source MariaDB.

\begin{tabularx}{\linewidth}{@{}Xrr@{}}
\toprule
Table PostgreSQL & MariaDB & Bronze \\
\midrule
\path{bronze.espocrm_account} & 1 & 1 \\
\path{bronze.espocrm_contact} & 2 & 2 \\
\path{bronze.espocrm_opportunity} & 1 & 1 \\
\bottomrule
\end{tabularx}

\preuve{crm/crm_01_tables_bronze.png}{Inventaire des tables CRM dans Bronze}
\quad
\preuve{crm/crm_02_volumes_bronze.png}{Comptage des lignes CRM dans Bronze}.

\section{Silver}

\subsection{Règles communes}

Cinq vues PostgreSQL versionnées sous \path{sql/silver/} exposent
une ligne par clé métier dans chaque entité. Elles calculent la
standardisation à partir de l'état courant de Bronze à chaque requête.
Après une synchronisation, les mêmes règles s'appliquent immédiatement
aux données chargées. Les colonnes suivent le \emph{snake\_case}.
Le système source, l'identifiant source et l'horodatage d'ingestion
assurent la traçabilité.

Dans les cinq vues, \path{ingested_at_utc} reprend directement
\path{_airbyte_extracted_at}, de type \code{timestamptz}. PostgreSQL
conserve l'instant et l'affiche dans le fuseau de la session. Le
paramètre SQL \verb|SET TIME ZONE 'UTC'| permet de l'afficher en UTC.
Les champs métier \path{created_at} et \path{modified_at} du CRM
gardent leur type source \code{timestamp without time zone}. Leur
conversion en instants UTC nécessite de confirmer le fuseau source.

Le dédoublonnage traite les occurrences répétées d'une même clé.
Une clé vide est écartée. Pour une clé présente plusieurs fois,
la vue retient l'extraction la plus récente et utilise l'identifiant
technique Airbyte pour départager les égalités.

\subsection{Flux 1 --- Vue des stations et prix des carburants}

\subsubsection{Grain et clé}

La vue \code{silver.api\_station}, définie dans
\path{sql/silver/api_station.sql}, reprend le modèle fourni.
Le grain est une ligne par station dans l'état courant chargé
en Bronze. \code{record\_id} et \code{source\_record\_id} reprennent
l'identifiant \code{id} de l'API en texte.

\subsubsection{Dictionnaire des colonnes et conversions}

{\small
\begin{tabularx}{\linewidth}{@{}>{\bfseries}p{3cm}X@{}}
\toprule
Groupe & Colonnes et règles appliquées \\
\midrule
Clés &
\code{record\_id} et \code{source\_record\_id} en \code{text},
issus de \code{id}. \\
Adresse &
\code{code\_postal}, \code{ville}, \code{adresse},
\code{code\_departement} et \code{code\_region} en \code{text}.
Le code postal provient de \code{cp}. \\
Prix &
\code{gazole\_prix}, \code{e10\_prix}, \code{e85\_prix},
\code{gplc\_prix}, \code{sp95\_prix} et \code{sp98\_prix}
en \code{numeric(12,3)}. \\
Horodatages métier &
Les six champs source \code{*\_maj} deviennent
\code{*\_mis\_a\_jour\_at} en \code{timestamptz}. \\
Disponibilité &
\code{carburants\_disponibles} et \code{carburants\_indisponibles}
en \code{jsonb}. \\
Traçabilité &
\code{source\_system} en \code{text},
\code{ingested\_at\_utc} en \code{timestamptz}
et \code{source\_file\_key} en \code{character varying}. \\
\bottomrule
\end{tabularx}
}

Les espaces en bordure des champs textuels sont retirés. Une chaîne
vide après nettoyage représente une valeur manquante et devient
\code{NULL}. Le type texte conserve les zéros initiaux, comme dans
le code postal \code{01000}. Le type \code{numeric(12,3)} conserve
la précision observée sur les prix, par exemple \code{2.301}.
Le décimal exact convient aux prix exprimés au millième. Le type
\code{double precision} utilise une représentation binaire approchée
qui peut introduire des écarts lors des comparaisons ou des calculs.
Un prix absent reste \code{NULL}.

L'horodatage \code{ingested\_at\_utc} reprend
\verb|_airbyte_extracted_at| en Bronze. Les horodatages métier
proviennent des champs de mise à jour de l'API. La valeur de
\code{source\_system} est \code{api\_prix\_carburants} et
\code{source\_file\_key} reprend la référence de l'objet MinIO.

La vue expose les champs nécessaires à l'identification des stations,
à leur localisation administrative et à l'analyse des prix et de
la disponibilité. Les autres champs de l'API, dont les horaires,
les services et le détail des ruptures, restent consultables en Bronze.

\subsubsection{Dédoublonnage et réconciliation}

La vue écarte les identifiants absents ou vides. Pour chaque
identifiant renseigné, \code{ROW\_NUMBER()} classe les lignes
par horodatage d'ingestion décroissant. L'identifiant technique
\verb|_airbyte_raw_id| départage les égalités. La vue conserve
la ligne de rang~1. Sur le lot contrôlé, aucune ligne n'a été écartée.

\subsection{Flux 2 --- Données du CSV}

La vue \path{silver.csv_departement} présente une ligne par code
département. Elle standardise les noms des sept champs métier en
\code{snake\_case}, supprime les espaces en bordure des valeurs
et convertit les chaînes vides en valeurs nulles. Les codes sont
conservés en texte pour préserver les zéros initiaux et les codes
alphanumériques.

Le code département fournit \path{record_id} et \path{source_record_id}.
La vue conserve également le système source, l'horodatage d'extraction
Airbyte et la référence du fichier MinIO. Les lignes sans code sont
écartées. Pour un même code, la vue conserve la ligne dont l'horodatage
d'extraction est le plus récent. En cas d'égalité, l'identifiant
technique Airbyte départage les lignes par ordre décroissant.
Le contrôle Bronze a relevé 101 codes uniques et aucune valeur absente.

Le script de création est \path{sql/silver/csv_departement.sql}.

\subsection{Flux 3 --- Entités CRM}

Les vues \path{silver.espocrm_account}, \path{silver.espocrm_contact}
et \path{silver.espocrm_opportunity} présentent une ligne par identifiant
source et par entité. Les champs textuels sélectionnés sont normalisés
par suppression des espaces en bordure et conversion des chaînes vides
en valeurs nulles. Les vues conservent le statut de suppression logique,
les clés de liaison disponibles et les champs de traçabilité du template.
Le montant conserve son type numérique et sa devise, et la date de
clôture conserve son type \code{date}. Les vues conservent les
clés de liaison en texte. Les contacts et les opportunités
référencent le compte par \path{account_id}, rapproché de
\path{record_id} dans la vue des comptes. Ce choix préserve le grain
de chaque entité et permet de vérifier séparément la présence
de la fiche et celle du compte attendu.

Les lignes sans identifiant sont écartées. Pour chaque identifiant, la
vue retient l'extraction Airbyte la plus récente et départage les
égalités par ordre décroissant de l'identifiant technique Airbyte.
Les trois scripts de création se trouvent dans \path{sql/silver/}.

\section{Tests}

\subsection{Flux 1 --- Contrôles du chargement Bronze}

{\small
\begin{tabularx}{\linewidth}{@{}Xrr@{}}
\toprule
Contrôle & Attendu & Observé \\
\midrule
Nombre de lignes dans Bronze & 9\,834 & 9\,834 \\
Identifiants de station distincts & 9\,834 & 9\,834 \\
Identifiants absents ou vides & 0 & 0 \\
Occurrences en doublon & 0 & 0 \\
Nombre de références de fichier distinctes & 1 & 1 \\
\bottomrule
\end{tabularx}
}

Les 9\,834 lignes référencent la clé MinIO
\path{raw/api/prix_carburants/extraction_date=2026-11-13/1794580622068_0.jsonl}.
Le comptage concorde avec celui du fichier téléchargé.

\preuve{api/api_08_carburants_controle_bronze.png}{Volume et clés de Bronze}. Les 9\,834 identifiants sont tous
présents et distincts.

\preuve{api/api_09_carburants_provenance_bronze.png}{Provenance du lot Bronze}. Les 9\,834 lignes référencent
un seul objet JSONL.

\begin{Verbatim}[fontsize=\small]
WITH stations AS (
    SELECT NULLIF(TRIM(_airbyte_data ->> 'id'), '') AS id_station
    FROM bronze.api_prix_carburants
)
SELECT COUNT(*) AS lignes_chargees,
       COUNT(DISTINCT id_station) AS ids_distincts,
       COUNT(*) FILTER (WHERE id_station IS NULL) AS ids_absents,
       COUNT(id_station) - COUNT(DISTINCT id_station)
           AS occurrences_en_doublon
FROM stations;

SELECT column_name, data_type
FROM information_schema.columns
WHERE table_schema = 'bronze'
  AND table_name = 'api_prix_carburants'
ORDER BY ordinal_position;

SELECT _ab_source_file_url AS fichier_source,
       COUNT(*) AS nombre_lignes
FROM bronze.api_prix_carburants
GROUP BY _ab_source_file_url;
\end{Verbatim}

\subsection{Flux 1 --- Contrôle préalable des conversions}

Le test a vérifié la conversion en \code{numeric(12,3)} de tous
les prix renseignés et la validité des horodatages renseignés.
Les valeurs absentes restent manquantes.

{\small
\begin{tabularx}{\linewidth}{@{}Xrrr@{}}
\toprule
Carburant & Prix renseignés & Prix invalides ou arrondis &
Horodatages invalides \\
\midrule
E10    & 7\,336 & 0 & 0 \\
E85    & 3\,762 & 0 & 0 \\
Gazole & 9\,319 & 0 & 0 \\
GPLc   & 1\,504 & 0 & 0 \\
SP95   & 2\,740 & 0 & 0 \\
SP98   & 7\,135 & 0 & 0 \\
\bottomrule
\end{tabularx}
}

\preuve{api/api_10_carburants_conversions.png}{Contrôle préalable des conversions}. Pour les six carburants,
la requête relève zéro prix invalide ou arrondi et zéro horodatage invalide.

\subsection{Flux 1 --- Résultats des huit contrôles Silver}

{\small
\begin{tabularx}{\linewidth}{@{}>{\bfseries}p{3cm}X@{}}
\toprule
Contrôle & Résultat observé et interprétation \\
\midrule
Clé non nulle & 0 clé absente en Silver. \\
Unicité & 9\,834 clés distinctes pour 9\,834 lignes. \\
Doublons & 0 occurrence en doublon en Bronze et en Silver sur ce lot. \\
Nommage & Les 24 noms sont conformes au \emph{snake\_case}. \\
Types & Les 24 types correspondent au dictionnaire Silver.
Les six prix ont une précision de~12 et une échelle de~3. \\
Réconciliation & 9\,834 lignes Bronze et 9\,834 lignes Silver.
Le test de réconciliation retourne vrai. \\
Traçabilité & 0 ligne sans traçabilité complète selon le contrôle
du système, de l'identifiant, de l'horodatage et du fichier source. \\
Règle métier & 0 station avec un prix renseigné négatif. \\
\bottomrule
\end{tabularx}
}

Les requêtes de bilan et d'inspection des colonnes figurent
dans \path{sql/tests/api_station.sql}. Les résultats portent sur
le lot contrôlé.

\preuve{api/api_11_carburants_controle_silver.png}{Bilan des contrôles Silver}. Le résultat confirme la
réconciliation des 9\,834 lignes, l'unicité des clés, la traçabilité
et l'absence de prix négatif.

\preuve{api/api_12_carburants_structure_silver.png}{Nommage et types des colonnes Silver}. L'inventaire confirme
24 noms conformes et le type \code{numeric(12,3)} des six prix.

\subsection{Flux 2 --- Contrôles de la vue Silver}

Le contrôle exécuté dans PostgreSQL retourne 101 lignes pour
101 identifiants distincts, comme en Bronze. Il relève zéro
identifiant absent, zéro occurrence en doublon et zéro ligne
sans les informations de traçabilité attendues. Les codes
\code{01}, \code{2A} et \code{2B} apparaissent chacun une fois.

Le contrôle est reproductible avec \path{sql/tests/csv_departement.sql}.
\preuve{csv/csv_05_departements_controle_silver.png}{Requête et résultat du contrôle Silver du CSV}.

\subsection{Flux 3 --- Contrôles des vues Silver}

La requête \path{sql/tests/espocrm.sql} confirme les volumes du lot
Bronze après le test du delta. Le nombre d'identifiants distincts
correspond au nombre de lignes pour chaque entité.

\begin{tabularx}{\linewidth}{@{}Xrrr@{}}
\toprule
Entité & Bronze & Silver & Identifiants distincts \\
\midrule
Comptes & 2 & 2 & 2 \\
Contacts & 4 & 4 & 4 \\
Opportunités & 2 & 2 & 2 \\
\bottomrule
\end{tabularx}

Pour chacune des trois vues, le contrôle relève zéro identifiant absent,
zéro doublon et zéro ligne présentant une anomalie sur le système source,
l'identifiant source ou l'horodatage d'ingestion.

\preuve{crm/crm_06_controles_silver.png}{Résultats des contrôles Silver du CRM}.

Le contrôle des rattachements, inclus dans le même script, retrouve
les deux contacts et l'opportunité créés pour le test du delta.
Pour les trois fiches, le compte référencé et le compte retrouvé
dans la vue Silver correspondent au compte de test
\code{6af78794aba6cf05e}.

\preuve{crm/crm_07_rattachements_silver.png}{Contrôle des trois rattachements au compte dans Silver}.

L'inventaire des trois vues confirme que les 37 noms de colonnes
respectent le format \code{snake\_case}. Les identifiants et clés de
liaison sont de type \code{text}. L'horodatage d'ingestion est de type
\code{timestamp with time zone}. Les champs de création et de
modification conservent leur type source \code{timestamp without time zone}.
Le montant de l'opportunité est de type \code{numeric}, la devise de
type \code{text}, la probabilité de type \code{bigint} et la date de
clôture de type \code{date}. Les indicateurs de suppression et
d'opposition aux appels sont de type \code{boolean}.

\preuve{crm/crm_08_types_nommage_silver.png}{Inventaire des noms et types des colonnes CRM dans Silver}.

\section{Delta CRM}

\subsection{État initial et nouvelles saisies}

Les relevés initiaux dans MariaDB et PostgreSQL Bronze montrent un
compte, deux contacts et une opportunité. Le test consiste à créer
un nouveau compte, deux contacts et une opportunité dans EspoCRM,
puis à rechercher leurs identifiants dans Bronze après synchronisation.

\subsection{Vérification de l'arrivée en Bronze}

\begin{tabularx}{\linewidth}{@{}Xrrr@{}}
\toprule
Entité & Avant & Après & Écart \\
\midrule
Comptes & 1 & 2 & +1 \\
Contacts & 2 & 4 & +2 \\
Opportunités & 1 & 2 & +1 \\
\bottomrule
\end{tabularx}

La requête de contrôle recherche les quatre identifiants EspoCRM
dans les tables correspondantes. Elle relève exactement une
occurrence par identifiant et confirme les écarts attendus.

\preuve{crm/crm_03_sync_delta_airbyte.png}{État des trois flux CRM dans Airbyte}.
L'écran indique deux comptes, quatre contacts et deux opportunités
chargés lors de la dernière synchronisation affichée.

\preuve{crm/crm_04_delta_identifiants_bronze.png}{Requête et résultat du contrôle des quatre identifiants CRM}.

Le contrôle de la relation entre l'opportunité et le compte joint
\path{bronze.espocrm_opportunity} à \path{bronze.espocrm_account}
sur \path{opportunity.account_id} et \path{account.id}. La requête
retourne une ligne et confirme que l'opportunité
\code{6af787ec38e9debdd} référence le compte de test
\code{6af78794aba6cf05e}.

\preuve{crm/crm_05_rattachement_opportunite_bronze.png}{Contrôle du rattachement de l'opportunité au compte dans Bronze}.

\section{Docker et Airbyte}
\label{sec:docker-airbyte}

\subsection{Installation locale et réseau}

Sur le poste Windows, Docker Desktop~4.92.0, Docker Engine~29.8.0
et Docker Compose~v5.5.1 ont été observés. Airbyte~2.5.0 est déployé
avec \code{abctl}. La stack a été démarrée depuis \path{compose/}
après création du fichier local \code{.env} à partir de
\code{.env.example}.

Airbyte est accessible sur le port~8000. MinIO expose son API sur
9000 et sa console sur 9001. PostgreSQL utilise 5432, MariaDB 3306,
EspoCRM 8081 et Adminer 8082. Depuis Airbyte sur Windows,
\code{host.docker.internal} permet de joindre les services exposés
sur la machine hôte. Le kit décrit les variantes macOS et Linux.

\href{https://github.com/MAHDIghr/projet_big_data_starter/blob/main/compose/README_DOCKER.md}{Procédure Docker du dépôt}

\subsection{Inventaire de l'environnement exécuté}

La validation de la configuration avec \code{docker compose config --quiet}
n'a affiché aucune erreur. Les six services Compose étaient démarrés.
PostgreSQL, MariaDB et EspoCRM affichaient l'état \code{healthy}.
Le contrôle \code{abctl local status} indiquait \code{deployed} pour
Airbyte, avec le chart Helm~2.5.0 et l'application~2.5.0.
Le contrôleur \code{ingress-nginx} utilisait le chart~4.15.1 et
l'application~1.15.1.

\subsubsection{Images Docker et version de l'outil de déploiement}

La sortie de \code{docker compose images} associe chaque conteneur
à l'image effectivement utilisée. Les colonnes sont resserrées
pour tenir sur la page.

\begin{Verbatim}[fontsize=\scriptsize]
CONTAINER            REPOSITORY       TAG      PLATFORM     IMAGE ID      SIZE     CREATED
dsid-adminer         adminer          6.1.0    linux/amd64  8df93d7bb431  164MB    7 weeks ago
dsid-espocrm         espocrm/espocrm  10.0.8   linux/amd64  817f2b063de2  1.08GB   2 months ago
dsid-espocrm-daemon  espocrm/espocrm  10.0.8   linux/amd64  817f2b063de2  1.08GB   2 months ago
dsid-mariadb         mariadb          11.8.9   linux/amd64  c4124d1afa87  453MB    2 weeks ago
dsid-minio           <none>           <none>   unknown      2d3964653720  0B       N/A
dsid-postgres        postgres         18.6     linux/amd64  e3d24f6906d0  644MB    8 days ago
\end{Verbatim}

La sortie de \code{abctl version} identifie l'outil utilisé pour
déployer Airbyte.

\begin{Verbatim}[fontsize=\small]
version: v0.30.4
\end{Verbatim}

\subsubsection{Identification de MinIO}

Le kit configure la construction depuis la référence
\path{RELEASE.2025-10-15T17-29-55Z}. L'image active \code{2d3964653720}
est affichée sans tag. La commande \code{docker exec dsid-minio minio --version}
identifie le binaire actif comme \code{DEVELOPMENT.GOGET}, compilé avec
Go~1.24.13 pour \code{linux/amd64}. Le tag de construction et
l'identification du binaire sont relevés séparément.

\subsection{Qualification de l'environnement de référence}

Les relevés Docker et Airbyte documentent la stack en fonctionnement
sur le poste Windows. Source S3~4.16.4 accède au lot API dans MinIO via
\url{http://host.docker.internal:9000}. Les contrôles SQL documentent
les chargements Bronze et les résultats Silver des trois sources.

\subsection{Incident de publication dans Builder}

\subsubsection{Diagnostic}

Le test Builder réussissait. L'action \emph{Publish} renvoyait
une erreur HTTP~500. Le journal serveur consulté avec
\code{kubectl logs} signalait l'absence de référence d'image
pour les manifestes de version majeure~7. Le YAML déclarait
la version~\code{7.28.2}.

\subsubsection{Correction et résultat}

\begin{enumerate}
  \item Vérifier que le pod PostgreSQL interne \code{airbyte-db-0}
        est en cours d'exécution.
  \item Vérifier l'empreinte SHA-256 de l'image
        \code{airbyte/source-declarative-manifest:7.28.2}
        avec \code{docker buildx imagetools inspect}.
  \item Ajouter la référence de la version majeure~7 dans la table
        \verb|declarative_manifest_image_version| de \code{db-airbyte},
        avec \code{ON CONFLICT DO NOTHING}.
  \item Relire l'entrée et relancer la publication.
\end{enumerate}

PowerShell a affiché \code{INSERT 0 1}, puis la version~\code{7.28.2}
associée à la version majeure~7. La publication a réussi avec
le même manifeste YAML.

\href{https://github.com/airbytehq/airbyte/issues/84928}{Ticket Airbyte et contournement documenté}

\subsection{Connexions et conventions de nommage}

\subsubsection{Connexion API vers MinIO}

La connexion affichée dans Airbyte est
\code{api\_prix\_carburants} \flux\ \code{minio\_raw}.
Elle relie \emph{API Prix des carburants France} à
\emph{MinIO RAW Prix des carburants}. Les 47 champs détectés
ont été sélectionnés. Le préfixe de flux \code{api\_} indique
l'origine des données.

Le déclenchement manuel permet de contrôler les exécutions.
Le mode \emph{Full Refresh -- Overwrite} relit le périmètre
national et maintient le dernier état dans la zone d'écriture.
Pour conserver un lot, archiver ses fichiers et son manifeste
hors de ce préfixe avant une nouvelle synchronisation.

\subsubsection{Chemin d'écriture et lecture S3}

JSONL conserve un objet par ligne et les structures imbriquées.
Le chemin de base est \path{raw/api/prix_carburants}. Le modèle
de chemin ajoute le dossier daté \path{extraction_date=YYYY-MM-DD}.
Le chemin de base identifie déjà le jeu, ce qui permet d'omettre
le nom du flux dans le modèle. Le job~5 a confirmé le chemin obtenu.

La lecture a été effectuée avec Source S3~4.16.4 et la destination
PostgreSQL~3.0.24. La source S3 utilise le bucket privé
\code{dsid-raw}, l'endpoint \url{http://host.docker.internal:9000}
et la région \code{us-east-1}. Elle réplique les enregistrements
JSONL. Les identifiants sont saisis dans les champs privés d'Airbyte.

\subsubsection{Noms techniques et noms affichés}

Les noms techniques utilisent des minuscules et des underscores.
Le modèle \verb|<origine>_to_<destination>| indique le sens du
transfert lorsqu'un nouveau nom technique est choisi. Le rapport
reprend les noms effectivement affichés dans Airbyte pour relier
les connexions aux journaux et aux preuves. Dans PostgreSQL,
le préfixe \code{api\_} identifie l'origine de la table Bronze.

\section{Reproductibilité}

\subsection{Ordre de reprise du flux API}

\begin{enumerate}
  \item Démarrer la stack du kit et vérifier les accès à Airbyte,
        MinIO et PostgreSQL.
  \item Importer le manifeste YAML dans Builder, tester le flux
        et publier le connecteur si nécessaire.
  \item Créer la source API et la connexion vers MinIO. Exécuter
        la collecte et relever le job et la clé de l'objet.
  \item Contrôler le volume et le SHA-256 du JSONL.
        Enregistrer les résultats dans le manifeste du lot.
  \item Exécuter la connexion Source S3 vers PostgreSQL Bronze.
        Conserver les trois champs entrants retenus et vérifier
        la structure, le volume, les clés et la provenance.
  \item Vérifier les conversions des six carburants, exécuter
        \path{sql/silver/api_station.sql}, puis les requêtes
        de \path{sql/tests/api_station.sql}.
\end{enumerate}

Le dépôt fournit le manifeste du connecteur sous
\path{airbyte/api_prix_des_carburants_france_v_2.yaml}.
Le motif S3 est réutilisable pour les collectes suivantes.
Pour rejouer exactement un lot, utiliser sa clé et vérifier son
SHA-256 à partir du manifeste. Le mode \emph{Overwrite} recharge
Bronze avec les objets sélectionnés. La vue Silver reflète ensuite
le contenu courant de Bronze.

\subsection{Vérification de la reprise et livrables}

Le
\href{https://github.com/MAHDIghr/projet_big_data_starter/blob/main/README.md}{guide de reprise}
décrit l'ordre de démarrage, la configuration des cinq connexions,
puis l'exécution des vues et contrôles SQL. Le
\href{https://github.com/MAHDIghr/projet_big_data_starter/blob/main/airbyte/README_AIRBYTE.md}{guide Airbyte}
précise les paramètres observés et ceux à vérifier.

Le démarrage et les trois chargements sont documentés sur le poste
Windows de référence. La vérification du dossier publié couvre les
chemins des guides, les scripts et les annexes. Les preuves disponibles
ne couvrent pas un rejeu intégral à partir d'une installation neuve.
La version exacte de la destination S3 n'est pas établie par ces relevés.

Pour le CSV, le guide décrit l'archivage de l'original, le contrôle
SHA-256 et le chargement de l'objet produit par Airbyte. Pour le CRM,
il décrit le contrôle initial, la création des fiches et la vérification
de leurs identifiants après synchronisation. Une nouvelle collecte
API produit un lot distinct dont le volume peut varier.

Le dossier technique accompagne le PDF. Il contient le README,
les fichiers Compose, \code{.env.example}, les configurations
Airbyte sans secrets, les scripts SQL, les manifestes et les preuves
associées aux résultats.

\end{document}